
\section{Discussion}
Three findings matter for anyone screening spontaneous reports for ototoxicity. First, adjusting for indication, ototoxic
co-medication, lawyer reports and demographics changed a substantial part of the hearing-loss signal list
(63 of 300 crude signals lost, 73 gained) but did not
improve agreement with labels. Second, replication in an independent national database behaved as a validation test should
for labelled associations---they replicated as often as their strength predicted---but unlabelled signals replicated less
than half as often as predicted, about as often as signals for drugs whose labels are silent on hearing. Third, replication is
not a clean validator on its own: solicited patient-support-programme reports produced ``replications'' for drugs with no
known association, and co-medication bias replicated whenever crude estimates were compared. After adjusted replication and
a review of report provenance, 4 of 22 robust unlabelled signals remained.

These results qualify the previous all-drug screens \cite{li2024,wu2025,liu2025} rather than contradict them. Their
labelled signals largely held; their unlabelled signals mostly did not, and several of the strongest are consistent with how
reports were collected rather than with pharmacology. Tafamidis and sacubitril/valsartan are examples: thousands of consumer
reports in FAERS, consistent with manufacturer programmes in the United States, and none in Canada. The same mechanism
produced spurious Canadian ``replications'' for biologics. Solicited reports are known to distort disproportionality
\cite{psp_solicited}, and their share differs between databases, so report provenance can manufacture both a signal and
its replication. Bias that is shared between databases---co-prescription in standard regimens---also replicates; etoposide
is the clearest example here.

Two inexpensive steps addressed most of what we found: comparing \emph{adjusted} estimates in both databases, and
profiling each case series for concentration---solicited or consumer reports, a single country or year, very drug-dense
reports. The power-calibrated comparison is also simple and avoids the outcome-dependent testability of rules based on
observed counts; it could be extended with empirical calibration against negative controls \cite{schuemie2014} or with
regression-based adjustment for co-reported drugs \cite{caster2010}. We suggest that screens for hearing outcomes report
these provenance indicators for every signal, in line with the case-by-case item of READUS-PV \cite{readuspv,readuspv_ee}.

The remaining candidates are hypotheses with plausible alternative explanations. Pregabalin, an $\alpha_2\delta$ ligand,
had 1{,}703 FAERS hearing-loss reports (adjusted OR
3.25 (3.09--3.43))
that replicated in Canada; a case of reversible, dose-related hearing loss has been described \cite{pregabalin_case}, but
diabetes, neuropathy and age confound \cite{diabetes_hearing}, and its labels list tinnitus and hyperacusis. For amlodipine,
whose labels mention tinnitus, a large cohort found a \emph{lower} risk of sensorineural hearing loss with calcium-channel
blockers \cite{ccb_cohort}, which argues against a causal effect. Doxycycline--tinnitus may reflect pulsatile tinnitus from
drug-induced intracranial hypertension, which is labelled \cite{doxy_iih}, or the infections it treats. None of these should
change practice without studies that have exposure data.

\paragraph{Limitations.} Spontaneous reports lack exposure denominators, and disproportionality measures reporting, not risk
or incidence. The label reference set is imperfect: labels are partly informed by spontaneous reports, list rare events
without causal assessment, and route-specific labels (for example topical gentamicin) blur ingredient-level status; we used
it only to compare methods. Excluding reports with an ototoxic co-medication conditions on a variable that reporters may add
after observing the event; stratifying on it instead gave the same conclusions. Indications are often missing, unreported
co-treatment (for example aminoglycosides in mycobacterial disease) remains, and content deduplication cannot catch
near-duplicate series, which the report-cluster indicator only partly addresses. The Canada Vigilance documentation does not
state that the extract is domestic; restricting it to reports sent directly to Health Canada gave similar results
(39/72 (54\%) of adjusted signals replicated, background
12/58 (21\%)). About 8{,}400 ingredients were screened without multiplicity
correction; the replication and calibration analyses, rather than individual $p$-values, carry the conclusions. Several
analyses---the solicited-report exclusion, the extended indication filter, adjusted-to-adjusted and calibrated replication,
the co-medication variants and the case-level review---were added after the primary results were seen, partly in response
to internal review; they are labelled post hoc and the pre-specified results are reported unchanged. The case-level review
used a written rubric applied by the author and by an AI assessor, not by independent clinicians.

\paragraph{Conclusions.} In FAERS, most drug--hearing-loss signals that are not already on a label do not survive adjustment,
adjusted replication in Canada Vigilance and a review of report provenance; disproportionality analysis remains a
hypothesis-generating tool. Replication in a second spontaneous-reporting database is informative only when both estimates
are adjusted and solicited reports are handled explicitly. The few remaining candidates, led by pregabalin, merit study with
exposure-based designs.
